\section{Abelian Varieties, the Theorem of the Cube, Isogenies, and Weil Bounds}
\label{sec:abelian-varieties-weil}

A proof of the Generalized Riemann Hypothesis built from unitary theta lifts and
PEL Shimura varieties cannot treat the geometry of abelian varieties or
resolution of singularities as informal background. Two separate pillars of the
global argument force the formalization to construct this geometry from the
ground up. On the arithmetic side, constructing Galois representations attached
to conjugate-self-dual automorphic forms requires abelian schemes over arbitrary
bases, polarizations, finite flat kernels of isogenies, $\ell$-adic Tate
modules, $p$-divisible groups, and characteristic polynomials of endomorphisms.
On the Arakelov-geometric side, controlling covolumes and arithmetic
intersection numbers on toroidal compactifications of Shimura varieties requires
functorial resolution of singularities in characteristic zero, Moriwaki's
arithmetic Hilbert--Samuel bigness theorem, and multiplicative height bounds of
Weil and Mahler type. Together these components account for a large share of the
reachable proof graph: $79{,}657$ modules in
\leanmod{Development.AbelianVarietyTheory}, $6{,}717$ in
\leanmod{Development.WeilBound}, $497$ in
\leanmod{Development.ResolutionOfSingularities}, $162$ in
\leanmod{GroundTheory.ArithmeticGeometry}, and $115$ in
\leanmod{ReciprocityTheory.SemistableReduction}.

We organize the algebraic-geometric architecture in four stages: spreading out
and the Theorem of the Cube, multiplication-by-$n$ and the isogeny calculus,
characteristic polynomials of endomorphisms and norm-degree identities, and
finally characteristic-zero resolution, multiplicative heights, and the
Jacquet--Shalika and Jacquet--Rallis inputs that bridge upward to
Section~\ref{sec:galois-cohomology-assembly}.

\subsection{Abelian schemes, spreading out, and the Theorem of the Cube}

Over a field $K$, an abelian variety $A$ of dimension $g$
\leanlink{GlobalTheory.AV.AbVar} is a smooth, proper, geometrically connected
group scheme $A \to \Spec K$ of relative dimension $g$. Properness and
geometrical connectedness force the group law on $A$ to be commutative. The
base case $g = 1$ of elliptic curves is developed explicitly in
\leanmod{GroundTheory.ArithmeticGeometry} via affine Weierstrass equations,
function-field transcendence of translates
\leanlink{GroundTheory.Arith.Abelian.transcendental_xOf_translate},
elliptic divisibility sequences and Somos recurrences for division points
\leanlink{GroundTheory.Arith.Abelian.uw_rel}
\leanlink{GroundTheory.Arith.Abelian.good}, and dual isogenies
\leanlink{GroundTheory.Arith.Abelian.mem_isogHom_iff}. For general dimension $g$
and over arbitrary base schemes $S$, the single highest-bottleneck theorem of
the entire algebraic-geometric layer is the Theorem of the Cube.

\begin{theorem}[Theorem of the Cube over an arbitrary base]
\leanlink{GlobalTheory.CubeTop.cubeTheorem}
\label{thm:cube-arbitrary-base}
Let $S$ be a scheme, let $A \to S$ be an abelian scheme, and let $\mathcal{L}$
be a line bundle on $A$. For every $S$-scheme $T$ and every triple of
$T$-valued points $a, b, c : T \to A$, the alternating tensor product
\[
  (a+b+c)^*\mathcal{L} \otimes (a+b)^*\mathcal{L}^{-1}
  \otimes (a+c)^*\mathcal{L}^{-1} \otimes (b+c)^*\mathcal{L}^{-1}
  \otimes a^*\mathcal{L} \otimes b^*\mathcal{L} \otimes c^*\mathcal{L}
  \otimes 0^*\mathcal{L}^{-1}
\]
is trivial on $T$ in the sense of the abelian-scheme line-bundle calculus
\leanlink{GlobalTheory.Published.CubeTheorem}.
\end{theorem}

\begin{proof}
The statement is established by reducing from an arbitrary base $S$ to
finitely generated models and then to perfect bases in
\leanlink{GlobalTheory.SpreadingModels.cubeTheorem_holds}.
\end{proof}

Why does the formalization place the spreading-out proof of
Theorem~\ref{thm:cube-arbitrary-base} inside \leanmod{Development.WeilBound}?
Because the reduction to finite-type models is a general scheme-theoretic
engine reused wherever rigidity or cohomology over general bases must be
descended from Noetherian or perfect rings.

\begin{proposition}[Spreading-out reduction for the Theorem of the Cube]
\leanlink{GlobalTheory.SpreadingModels.cubeTheorem_holds}
\label{prop:cube-spreading}
The universal Theorem of the Cube \leanlink{GlobalTheory.Published.CubeTheorem}
holds over every base scheme $S$.
\end{proposition}

\begin{proof}
One first spreads out the scheme data \leanlink{GlobalTheory.T.scheme_models_spread},
smoothness \leanlink{GlobalTheory.T.smooth_spreads_out}, properness
\leanlink{GlobalTheory.T.isProper_spreads_out}, and the line bundle
\leanlink{GlobalTheory.T.line_bundle_spreads_out'} from an arbitrary base to
finitely presented models \leanlink{GlobalTheory.T.has_model_of_spreading}.
From finite-type models one deduces the statement over perfect bases
\leanlink{GlobalTheory.T.perfect_of_models}, and finally passes from perfect
bases to arbitrary bases \leanlink{GlobalTheory.T.cube_theorem_of_perfect}.
\end{proof}

The Theorem of the Cube drives the entire Picard and intersection calculus of
abelian varieties. Given a class $c \in \Pic(A)$ and two $Y$-points
$x, y : Y \to A$, the Mumford class
\[
  \Lambda(c; x, y) := (x + y)^* c - x^* c - y^* c \in \Pic(Y)
\]
measures the failure of pullback along $c$ to be additive. By
Theorem~\ref{thm:cube-arbitrary-base}, $\Lambda(c; x, y)$ is biadditive in $x$
and $y$.

\begin{lemma}[Additivity of the Mumford class]
\leanlink{GlobalTheory.AV.AbVar.lam_mul_left_of_cube}
\label{lem:mumford-additive}
Let $A$ be an abelian variety over a field $K$ and let $c \in \Pic(A)$. For any
$K$-scheme $Y$ and any $Y$-points $x, x', y : Y \to A$,
\[
  \Lambda(c; x + x', y) = \Lambda(c; x, y) + \Lambda(c; x', y).
\]
\end{lemma}

\begin{proof}
Lift $c \in \Pic(A)$ to a line bundle $\mathcal{L}$ via
\leanlink{GlobalTheory.CechPic.cls_surj} and apply
\leanlink{GlobalTheory.CubeTop.cubeTheorem} to the triple $(x, x', y)$. Passing
back to $\Pic(Y)$ via \leanlink{GlobalTheory.CechPic.cls_linComb} yields the
eight-term alternating identity, which rearranges directly into
$\Lambda(c; x + x', y) - \Lambda(c; x, y) - \Lambda(c; x', y) = 0$.
\end{proof}

A symmetric line bundle class $d \in \Pic(B)$, meaning $[-1]_B^* d = d$,
therefore satisfies the quadratic pullback law
\[
  (f + g)^* d + (f - g)^* d = 2 f^* d + 2 g^* d
\]
and $(n f)^* d = n^2 f^* d$ for all $f, g \in \Hom(A, B)$ and $n \in \Z$
\leanlink{GlobalTheory.T.mumfordClass_symmetric_pullback_quadratic}. To turn
this algebraic identity into a metric on $\Hom(A, B)$, one first needs
projectivity and an ample symmetric class. Over an algebraically closed field,
projectivity of every abelian scheme is proved directly via charts in
\leanlink{GlobalTheory.AbelianProjective.isProjBy_of_isAbelianScheme}, and then
descended to an arbitrary base field $K$ using norm-chart base change
\leanlink{GlobalTheory.ProjField.isProjBy_of_isAbelianScheme}. Symmetrizing an
ample class $c_0 \in \Pic(A)$ via $c := c_0 + [-1]_A^* c_0$ produces a
symmetric ample class on every abelian variety.

\begin{proposition}[Existence of symmetric ample classes]
\leanlink{GlobalTheory.AV.AbVar.exists_symm_ample}
\label{prop:symm-ample}
For every field $K$ and every abelian variety $A$ over $K$, there exists an
ample class $c \in \Pic(A)$ satisfying $[-1]_A^* c = c$.
\end{proposition}

\begin{proof}
Choose an ample class $c_0 \in \Pic(A)$ from
\leanlink{GlobalTheory.T.abelianVariety_exists_ampleClass}. Since
$[-1]_A \circ [-1]_A = \mathrm{id}_A$
\leanlink{GlobalTheory.AV.AbVar.neg_one_left_comp_self}, the involution $[-1]_A$
is a finite automorphism, so $[-1]_A^* c_0$ is again ample and the sum
$c = c_0 + [-1]_A^* c_0$ is ample and fixed by $[-1]_A^*$ by
\leanlink{GlobalTheory.CechPic.pull_comp} and
\leanlink{GlobalTheory.CechPic.pull_id}.
\end{proof}

Combining a symmetric ample class $d \in \Pic(B)$ with an ample class
$c \in \Pic(A)$ and the multilinear symmetric intersection number on $A$
\leanlink{GlobalTheory.AV.AbVar.inter_update_add} produces a positive-definite
integer-valued quadratic form $Q : \Hom(A, B) \to \Z$
\leanlink{GlobalTheory.AV.AbVar.exists_posDef_quadratic_hom} (exported as
\leanlink{GlobalTheory.T.exists_posDef_quadraticForm_hom}). Explicitly, when
$\dim A = 0$ every homomorphism vanishes and $Q \equiv 0$; when
$\dim A = n + 1 \ge 1$, one takes a symmetric ample class $d \in \Pic(B)$ from
Proposition~\ref{prop:symm-ample} and an ample class $c \in \Pic(A)$
\leanlink{GlobalTheory.AV.exists_isAmpleCls}, and sets
\[
  Q(f) := \bigl(f^* d \cdot c^n\bigr)_A.
\]
Multilinearity of intersection numbers together with the quadratic pullback
identity \leanlink{GlobalTheory.AV.AbVar.pull_add_add_pull_sub} yields the
parallelogram law
\[
  Q(f + g) + Q(f - g) = 2 Q(f) + 2 Q(g)
\]
for all $f, g \in \Hom(A, B)$, while
\leanlink{GlobalTheory.T.mumfordClass_symmetric_pullback_quadratic} gives
$Q(0) = 0$ and positivity of ample intersection degrees gives $Q(f) > 0$ for
every $f \neq 0$. Combined with torsion-freeness of $\Hom(A, B)$
\leanlink{GlobalTheory.AV.AbVar.hom_torsionFree}, this positive-definite
quadratic form implies that $\Hom(A, B)$ is a free $\Z$-module of finite rank
at most $4 (\dim A)(\dim B)$
\leanlink{GlobalTheory.T.hom_finite_free_finrank_le_four_mul_dim}. That finite
rank is later essential for proving integrality of characteristic polynomials
of endomorphisms.

\subsection{Isogenies, multiplication-by-$n$, torsion, and $p$-divisible groups}

A homomorphism $f : A \to B$ of abelian varieties over $K$ is an
\emph{isogeny} \leanlink{GlobalTheory.AV.AbVar.IsIsogeny} if the underlying
scheme morphism is finite and surjective. Its degree $\deg(f)$
\leanlink{GlobalTheory.AV.AbVar.deg} is the generic-fiber rank of the finite
flat morphism $f$. The cornerstone calculation of the isogeny theory is the
degree of multiplication by a nonzero integer.

\begin{theorem}[Multiplication-by-$n$ is an isogeny of degree $|n|^{2g}$]
\leanlink{GlobalTheory.MulIsogeny.mul_isogeny}
\label{thm:mul-isogeny}
Let $A$ be an abelian variety of dimension $g$ over a field $K$. For every
integer $n \in \Z$ with $n \neq 0$, the endomorphism $[n]_A := n \cdot \mathrm{id}_A$
is an isogeny of degree
\[
  \deg([n]_A) = |n|^{2g}.
\]
\end{theorem}

\begin{proof}
For $n > 0$, one base-changes to an algebraic closure $\overline{K}$ using
isogeny and degree invariance under field extension
\leanlink{GlobalTheory.AV.AbVar.isIsogeny_bcMap_iff_and_deg_bcMap}. Over
$\overline{K}$, multiplication by $n$ is finite and flat, hence surjective
because $A_{\overline{K}}$ is geometrically integral, and its generic-fiber
rank equals $n^{2g}$. For $n < 0$, write $[n]_A = [-1]_A \circ [|n|]_A$. The
inversion $[-1]_A$ is an involution and therefore an isogeny of degree $1$ by
\leanlink{GlobalTheory.MulIsogeny.neg_one}, so multiplicativity of isogeny
degrees gives $\deg([n]_A) = |n|^{2g}$.
\end{proof}

When $n$ is invertible in $K$, the finite flat group scheme $A[n] = \ker([n]_A)$
is étale, so its geometric point count equals the degree of $[n]_A$.

\begin{theorem}[Geometric $n$-torsion cardinality]
\leanlink{GlobalTheory.AV.AbVar.card_torsion_pt}
\label{thm:card-torsion}
Let $A$ be an abelian variety of dimension $g$ over a field $K$. For every
positive integer $n$ with $n \neq 0$ in $K$, the group of geometric $n$-torsion
points satisfies
\[
  \bigl|A[n](\overline{K})\bigr| = n^{2g}.
\]
\end{theorem}

\begin{proof}
Since $n \neq 0$ in $K$, we have $n \neq 0$ in $\Z$, so $[n]_A$ is an isogeny
of degree $n^{2g}$. The action of $[n]_A$ on points is raising to the $n$-th
power by \leanlink{GlobalTheory.AV.AbVar.comp_pow_id_torscard}, and the kernel
is killed by the invertible integer $n$, so the degree of $[n]_A$ coincides
with the cardinality of its geometric kernel.
\end{proof}

Two structural consequences follow immediately. First, any homomorphism
$f : A \to C$ vanishing on $A[n](\overline{K})$ for an invertible integer $n$
is divisible by $n$ in $\Hom(A, C)$
\leanlink{GlobalTheory.AV.AbVar.exists_nsmul_eq_of_kills_torsion}: because
$A[n]$ is reduced, vanishing on geometric $n$-torsion forces vanishing on the
scheme-theoretic kernel of $[n]_A$, allowing $f$ to factor through the isogeny
$[n]_A$. Second, every isogeny $f : A \to B$ admits a quasi-inverse
\leanlink{GlobalTheory.AV.AbVar.exists_quasiInverse_isogeny} (exported as
\leanlink{GlobalTheory.T.exists_quasiInverse_isogeny'}): there exist an
integer $n > 0$ and an isogeny $g : B \to A$ satisfying
\[
  g \circ f = [n]_A \qquad \text{and} \qquad f \circ g = [n]_B.
\]
To construct $g$, one observes that $A$ and $B$ have equal dimension and that
the finite kernel of $f$ is annihilated by some integer $n > 0$
\leanlink{GlobalTheory.AV.KerKill.kernel_killed}. Flat-surjective descent over
$\Spec K$ in \leanlink{GlobalTheory.IsogenyQuasiInverse.exists_quasiInverse}
then produces a homomorphism $g : B \to A$ with $g \circ f = [n]_A$ and
$f \circ g = [n]_B$ (verified via \leanlink{GlobalTheory.AV.AbVar.hom_ext});
since $[n]_B = f \circ g$ is an isogeny by Theorem~\ref{thm:mul-isogeny}, the
factor $g$ is also an isogeny.

Theorem~\ref{thm:card-torsion} and the existence of quasi-inverse isogenies
connect the scheme-theoretic category of abelian varieties to $\ell$-adic
linear algebra and $p$-adic Dieudonné theory. For any prime $\ell$ invertible
in $K$, the $\ell$-adic Tate module
$T_\ell(A) = \lim_{\leftarrow r} A[\ell^r](\overline{K})$ is a free
$\Z_\ell$-module of rank $2g$, and
$V_\ell(A) = T_\ell(A) \otimes_{\Z_\ell} \Q_\ell$ is a $2g$-dimensional
$\Q_\ell$-vector space.

\begin{proposition}[Tate-module criterion for finiteness and isogenies]
\leanlink{GlobalTheory.AV.AbVar.isFinite_iff_tmap_injective_and_isIsogeny_iff_vmap_bijective}
\label{prop:tate-isogeny}
Let $K$ be a field, let $\ell$ be a prime with $\ell \neq 0$ in $K$, and let
$f : A \to B$ be a homomorphism of abelian varieties over $K$. Then:
\begin{enumerate}
  \item $f$ is finite if and only if the induced map
  $T_\ell(f) : T_\ell(A) \to T_\ell(B)$ is injective;
  \item $f$ is an isogeny if and only if
  $V_\ell(f) : V_\ell(A) \to V_\ell(B)$ is bijective.
\end{enumerate}
\end{proposition}

\begin{proof}
If $f$ is an isogeny, its quasi-inverse $g : B \to A$ with $g \circ f = [n]_A$
and $f \circ g = [n]_B$ implies that $V_\ell(f)$ is bijective because
multiplication by $n > 0$ is invertible on $V_\ell(A)$ and $V_\ell(B)$
\leanlink{GlobalTheory.AV.AbVar.IsogTate.vmap_bijective_of_isIsogeny}.
Conversely, if $V_\ell(f)$ is bijective, then $T_\ell(f)$ is injective by
\leanlink{GlobalTheory.AV.AbVar.IsogTate.tmap_injective_of_vmap_injective},
making $f$ finite, and comparing ranks $2 \dim A = \dim_{\Q_\ell} V_\ell(A) =
\dim_{\Q_\ell} V_\ell(B) = 2 \dim B$ forces $\dim A = \dim B$, hence
surjectivity of $f$.
\end{proof}

At the residue characteristic $p$ (where $[p^r]_A$ is no longer étale), the
role of $T_\ell(A)$ is played by the $p$-divisible group $A[p^\infty]$.

\begin{proposition}[The $p$-divisible group of an abelian variety]
\leanlink{GlobalTheory.AVTower.exists_tower}
\label{prop:pdiv-tower}
For every prime $p$, every field $k$, and every abelian variety $A$ of
dimension $g$ over $k$, there exists a $p$-divisible group $G$ over $k$ of
height $2g$ realized as a torsion tower of $A$ (representing the system
$\{A[p^n]\}_{n \ge 1}$).
\end{proposition}

\begin{proof}
At each level $n$, the finite flat kernel $A[p^n]$ has $k$-rank $p^{2gn}$ by
\leanlink{GlobalTheory.AVTower.finrank_torsion}, proved from
Theorem~\ref{thm:mul-isogeny}. The compatible transition maps yield a
$p$-divisible group $G$ whose level-$1$ rank $p^{\mathrm{ht}(G)}$ equals
$p^{2g}$, fixing $\mathrm{ht}(G) = 2g$.
\end{proof}

\subsection{Characteristic polynomials of endomorphisms and norm-degree identities}

On PEL Shimura varieties, Frobenius endomorphisms and Hecke correspondences
act on the abelian varieties parameterized by moduli points, and one must
compute their characteristic polynomials independently of $\ell$ and relate the
degrees of endomorphisms coming from an order $\mathcal{O}_F$ to field norms.
In Lean, a polynomial $P \in \Z[X]$ is called a characteristic polynomial of
$f \in \End(A)$ \leanlink{GlobalTheory.AV.AbVar.IsCharpoly} if $P$ is monic of
degree $2g$ and satisfies
\[
  P(n) = \deg([n]_A - f) \qquad \text{for all } n \in \Z.
\]

\begin{theorem}[Existence and uniqueness of the characteristic polynomial]
\leanlink{GlobalTheory.AV.AbVar.isCharpoly_charpoly_unique}
\label{thm:charpoly-unique}
For every field $K$, every abelian variety $A$ of dimension $g$ over $K$, and
every endomorphism $f \in \End(A)$, there exists a unique monic polynomial
$P_f \in \Z[X]$ of degree $2g$ satisfying $P_f(n) = \deg([n]_A - f)$ for all
$n \in \Z$, namely $P_f = \mathrm{charpoly}(f)$.
\end{theorem}

\begin{proof}
Existence is proved in \leanlink{GlobalTheory.AV.AbVar.exists_isCharpoly} by
combining three ingredients. First, the degree of an integer linear combination
of homomorphisms is a homogeneous polynomial of degree $2g$
\leanlink{GlobalTheory.T.homDegree_integerCombination_isHomogeneous}, which
yields a monic rational polynomial $P \in \Q[X]$ of degree $2g$ interpolating
$\deg([n]_A - f)$. Second, one picks an auxiliary prime $\ell$ invertible in
$K$ (taking $\ell = 2$ if $\operatorname{char} K \neq 2$ and $\ell = 3$
otherwise) and identifies $P$ over $\Q_\ell$ with the characteristic polynomial
of $V_\ell(f)$ on the $2g$-dimensional $\Q_\ell$-space $V_\ell(A)$. Third,
because $\End(A)$ is a finitely generated free $\Z$-module, $f$ satisfies a
monic integral polynomial $Q \in \Z[X]$, which annihilates $V_\ell(f)$ and forces
every coefficient of $P$ to be an integer. Uniqueness then follows immediately
because two polynomials over $\Z$ agreeing at all $n \in \Z$ are identical.
\end{proof}

The published interface to Theorem~\ref{thm:charpoly-unique} is exported as
\leanlink{GlobalTheory.T.abelianVarietyEndCharpoly_spec}. A clean application
is the norm-degree identity for abelian varieties with real or complex
multiplication by the ring of integers $\mathcal{O}_F$ of a number field $F$.

\begin{theorem}[Norm-degree identity for number-field actions]
\leanlink{GlobalTheory.AV.AbVar.DegNorm.deg_pow_finrank_eq_natAbs_norm_pow}
\label{thm:norm-degree}
Let $B$ be an abelian variety of dimension $g$ over a field $K$, let $F$ be a
number field, and let $\psi : \mathcal{O}_F \to \End(B)$ be a ring
homomorphism. For every nonzero element $c \in \mathcal{O}_F$,
\[
  \deg(\psi(c))^{[F:\Q]} = \bigl|N_{F/\Q}(c)\bigr|^{2g}.
\]
\end{theorem}

\begin{proof}
When $g = 0$, choosing $y \in \mathcal{O}_F$ with $c y = N_{F/\Q}(c)$ gives
$\psi(y) \circ \psi(c) = [N_{F/\Q}(c)]_B$, which has degree $1$ by
\leanlink{GlobalTheory.MulIsogeny.mul_isogeny}, forcing $\deg(\psi(c)) = 1$.
When $g > 0$, the characteristic polynomial of $\psi(c)$ from
\leanlink{GlobalTheory.AV.AbVar.isCharpoly_charpoly_unique} is a power
$m_c^k$ of the minimal polynomial $m_c$ of $c$ over $\Q$, with
$k [\Q(c):\Q] = 2g$. Equating $\deg(\psi(c))$ with the product of roots of
$P_{\psi(c)}$ in $\overline{\Q}$ and comparing with the norm formula
$N_{F/\Q}(c) = N_{\Q(c)/\Q}(c)^{[F:\Q(c)]}$ yields
$\deg(\psi(c))^{[F:\Q(c)]} = |N_{F/\Q}(c)|^k$; raising both sides to the power
$[\Q(c):\Q]$ and substituting $k [\Q(c):\Q] = 2g$ gives the result.
\end{proof}

\subsection{Resolution of singularities, multiplicative heights, and Arakelov--Weil bounds}

The algebraic geometry of abelian varieties handles the moduli objects
themselves; the remaining three domains of this chapter supply the compactification,
volume, and automorphic matching tools needed to extract quantitative bounds
from them.

\paragraph{Functorial resolution of singularities in characteristic zero.}
In \leanmod{Development.ResolutionOfSingularities}, the proof builds a
scheme-theoretic blowup and embedded-resolution engine from scratch. The
foundational existence theorem constructs the blowup along any ideal sheaf by
gluing local $\operatorname{Proj}$ schemes.

\begin{theorem}[Existence of blowups along ideal sheaves]
\leanlink{FunctorialResolution.Blowup.exists_isBlowupAlongIdealSheaf}
\label{thm:blowup-exists}
For every scheme $X$ and every ideal sheaf $\mathcal{I}$ on $X$, there exists a
morphism of schemes $\pi : Y \to X$ that is a blowup of $X$ along $\mathcal{I}$.
\end{theorem}

\begin{proof}
The scheme $Y$ and projection $\pi$ are constructed by gluing the local
projective blowups over affine opens, verified in
\leanlink{FunctorialResolution.Blowup.isBlowupAlongIdealSheaf_gluedProj}.
\end{proof}

Building on blowup existence, étale coordinate charts
\leanlink{FunctorialResolution.EmbeddedResolution.isSmoothDivisor_of_locally_coordIdeal},
monomial boundary order analysis
\leanlink{FunctorialResolution.EmbeddedResolution.MonomialChart.orderAt_boundary_eq_one},
and canonical divisor reduction
\leanlink{FunctorialResolution.EmbeddedResolution.divisorRestrictionFieldBaseChange},
the proof obtains characteristic-zero resolution of singularities and
principalization of ideal sheaves in
\leanlink{PositivityTheory.ResPlug.FromRes.resK}.

\paragraph{Moriwaki's arithmetic bigness theorem and Arakelov volumes.}
Resolution of singularities in characteristic zero plugs directly into the
Arakelov-geometric heart of \leanmod{Development.WeilBound}. Pairing
\leanlink{PositivityTheory.ResPlug.FromRes.resK} with the Poincaré--Lelong
arithmetic intersection system yields Moriwaki's arithmetic bigness theorem
and the accompanying arithmetic Hilbert--Samuel lower bound
\leanlink{PositivityTheory.Pf1.HilbertSamuelMoriwaki.hsUse_outright}.

\begin{theorem}[Moriwaki arithmetic bigness]
\leanlink{PositivityTheory.Pf1.HilbertSamuelMoriwaki.moriwakiBig_outright}
\label{thm:moriwaki-big}
The arithmetic bigness criterion of Moriwaki
(\leanlink{ReciprocityTheory.ChartRestriction.G.ArakelovGeometry.MoriwakiBigStatement})
holds unconditionally for Hermitian line bundles on arithmetic varieties.
\end{theorem}

\begin{proof}
The proof combines characteristic-zero resolution of singularities
\leanlink{PositivityTheory.ResPlug.FromRes.resK} with the arithmetic
intersection system derived from the Poincaré--Lelong equation and the dense
curve criterion in
\leanlink{PositivityTheory.Pf1.HilbertSamuelMoriwaki.moriwakiBig_of_classicalTheorems}.
\end{proof}

In the global architecture, Theorem~\ref{thm:moriwaki-big} and
\leanlink{PositivityTheory.Pf1.HilbertSamuelMoriwaki.hsUse_outright} feed
directly into the self-intersection volume estimates on toroidal
compactifications of Siegel and unitary Shimura varieties inside
\leanlink{UpperTheory.LiteratureTheorems.every_character_of_literature_of_vWin},
supplying the lower bounds on arithmetic volumes that balance the Bost slope
inequality.

\paragraph{Multiplicative heights, Jacquet--Shalika rigidity, and Jacquet--Rallis transfer.}
Three auxiliary pillars in this layer feed the automorphic and archimedean
estimates:
\begin{itemize}
  \item \emph{Multiplicative heights and Borel--Weil polynomials.} The package
  \leanmod{ReciprocityTheory.SemistableReduction} formalizes Hermitian forms
  \leanlink{ReciprocityTheory.MultiplicativeHeights.HermForm}, Gauss absolute
  values \leanlink{ReciprocityTheory.MultiplicativeHeights.gaussAbs}, Mahler
  torus means \leanlink{ReciprocityTheory.MultiplicativeHeights.torusMean},
  and block-sphere logarithmic averages
  \leanlink{ReciprocityTheory.MultiplicativeHeights.blockSphereMean} for
  multihomogeneous polynomials
  \leanlink{ReciprocityTheory.MultiplicativeHeights.IsMultiHomogeneous}
  (together with almost-everywhere nonvanishing of nonzero multivariate
  polynomials
  \leanlink{ReciprocityTheory.MultiplicativeHeights.ae_eval_ne_zero_of_ne_zero_fin}).
  \item \emph{Jacquet--Shalika rigidity and Rankin--Selberg poles.} Inside
  \leanmod{Development.WeilBound}, the proof establishes Rankin--Selberg pole
  detection \leanlink{PositivityTheory.AdelicQuotientData.rsPole} and
  Jacquet--Shalika strong multiplicity-one rigidity
  \leanlink{PositivityTheory.AdelicQuotientData.jacquetShalikaRigidity}, which
  guarantee that cuspidal automorphic representations of $\GL_n$ are uniquely
  pinned by almost all local spherical parameters.
  \item \emph{Jacquet--Rallis unit matching and rank-one transfer.} Within
  \leanmod{Development.AbelianVarietyTheory.Proofs.InertiaEigenvalues}, the
  fundamental lemma and unit matching for the Jacquet--Rallis relative trace
  formula are established in all ranks \leanlink{GlobalTheory.JR.unitMatch_all}
  (exported as \leanlink{GlobalTheory.T.jacquetRallisUnitMatch_all_ranks})
  from the quadratic-field Weil representation
  \leanlink{GlobalTheory.T.jacquetRallisUnitMatch_of_weilRep_quadratic},
  together with rank-one étale orbital Fourier transfer
  \leanlink{GlobalTheory.JR.RankOneTransfer.rankOne}.
\end{itemize}

Table~\ref{tab:ch04-pillars} summarizes how the main algebraic-geometric
theorems of this chapter feed upward into the Galois-theoretic and
cohomological assembly of Section~\ref{sec:galois-cohomology-assembly}.

\begin{table}[ht]
\centering
\caption{Algebraic-geometric pillars of Section~\ref{sec:abelian-varieties-weil} and their downstream roles.}
\label{tab:ch04-pillars}
{\small
\begin{tabular}{@{}p{0.23\textwidth}>{\raggedright\arraybackslash}p{0.41\textwidth}>{\raggedright\arraybackslash}p{0.30\textwidth}@{}}
\toprule
Result & Lean anchor & Downstream consumer \\
\midrule
Theorem of the Cube & \leanlink{GlobalTheory.CubeTop.cubeTheorem} & Mumford biextensions, polarizations \\
Spreading out & \leanlink{GlobalTheory.SpreadingModels.cubeTheorem_holds} & Arbitrary-base Cube Theorem \\
Degree of $[n]_A$ & \leanlink{GlobalTheory.MulIsogeny.mul_isogeny} & Torsion counts, $p$-divisible height \\
$n$-torsion points & \leanlink{GlobalTheory.AV.AbVar.card_torsion_pt} & $\ell$-adic Tate modules, level structures \\
Quasi-inverse isogeny & \leanlink{GlobalTheory.AV.AbVar.exists_quasiInverse_isogeny} & $V_\ell$ bijectivity, Dieudonné integrality \\
Char.\ polynomial & \leanlink{GlobalTheory.AV.AbVar.isCharpoly_charpoly_unique} & Norm-degree law, Frobenius eigenvalues \\
Blowups \& resolution & \leanlink{PositivityTheory.ResPlug.FromRes.resK} & Moriwaki arithmetic bigness \\
Moriwaki bigness & \leanlink{PositivityTheory.Pf1.HilbertSamuelMoriwaki.moriwakiBig_outright} & Shimura-variety arithmetic volumes \\
\bottomrule
\end{tabular}}
\end{table}
